% ============================================================================
\clearpage
\section*{v4.9.0-r30 Source-Only: The Compensator Avoids the Classical One-Half, and Why the Exchange Is Coherent}
\addcontentsline{toc}{section}{v4.9.0-r30 Source-Only: The Compensator Avoids the Classical One-Half, and Why the Exchange Is Coherent}
\label{sec:v490-r30-compensator-coherence}
% ============================================================================

\paragraph{Purpose.}
r29 left open why the exchange between matter and tachyon should be coherent rather than a classical replacement.  On 2026-10-22
the author pointed out that the Heisenberg compensator was introduced precisely to avoid the classical probability one half of
finding the object in one or the other sector.  This layer makes that precise, checks it against the document's own use of the
compensator, and draws the consequence for r28 and r29.  The script is \texttt{llm/k67\_kompensator\_kohaerenz.py}.  No new
prediction is made, and nothing is compared with data.

\noindent\textbf{Source-only status.}
This v4.9.0-r30 layer is source-only.  It does not build a PDF, does not assign a new DOI,
and does not authorize a public release.  The v4.9.0 Version DOI
\href{https://doi.org/10.5281/zenodo.23534954}{10.5281/zenodo.23534954} and the Concept DOI
\href{https://doi.org/10.5281/zenodo.19210728}{10.5281/zenodo.19210728} are unchanged by it.

\subsection*{v4.9.0-r30 classical one-half against the compensator}
\label{subsec:v490-r30-one-half}

\paragraph{Theoremlet v4.9.0-r30.1.}
\begin{enumerate}[leftmargin=2em]
\item[(a)] A classical mixture of a state with its tachyon image, \(\tfrac12\rho+\tfrac12\sigma_y\rho\sigma_y\), puts weight \(1/2\) into the
compensator port for every state.  This is the classical one half.
\item[(b)] A coherent combination of the amplitudes \(f\) and \(e^{i\varphi}f_T\) puts a weight into the compensator port that depends on the
relative phase.  It is \(1/2\) only at two phases, and it is \(0\) or \(1\), a definite outcome, at the two others.
\item[(c)] The document's own use of the compensator is coherent.  The remark ``Light as the Heisenberg compensator of the universe''
of the main body computes \(C(\bsym)=\tfrac12(\bsym-\bsym)=0\).  The cancellation exists only for amplitudes.  A mixture of
probabilities gives \(\tfrac12\bsym^2+\tfrac12\bsym^2=\bsym^2=0.302\), never \(0\).
\end{enumerate}

\paragraph{Proof sketch.}
(a) \(\sigma_yC_{\HC}\sigma_y=A\), so the mixture has weight \(\tfrac12(\operatorname{tr}C_{\HC}\rho+\operatorname{tr}A\rho)=\tfrac12\); checked on \(1000\)
random states.  (b) numerical at \(\varphi=0,\pi/2,\pi,3\pi/2\): \(1/2,0,1/2,1\).  (c) direct.  \(\square\)

\paragraph{Reading.}
This confirms the author's account.  The compensator acts on amplitudes.  Where a classical description can only say ``with
probability one half here or there'', the compensator gives a definite port, decided by the relative phase.  The remark on light
uses exactly this property.

\subsection*{v4.9.0-r30 the compensator partner is the tachyon seen through time reversal}
\label{subsec:v490-r30-partner}

\paragraph{Theoremlet v4.9.0-r30.2.}
\begin{enumerate}[leftmargin=2em]
\item[(a)] \(C=\tfrac12(1+J)\) is a projection only for an involution, \(J^2=1\).  This holds for \(J_B=\sigma_x\).  It fails for the tachyon map
\(z\mapsto iz\), whose square is \(-1\) on amplitudes; there \(C^2=J/2\).
\item[(b)] \(J_B\) is the tachyon map followed by time reversal: \(J_B=R_{\pi/2}\circ T\), with \(T\colon z\mapsto\bar z\).  On the doublet,
\(\left(\begin{smallmatrix}0&-1\\1&0\end{smallmatrix}\right)\operatorname{diag}(1,-1)=\sigma_x\).
\item[(c)] So the compensator combines a state with its tachyon partner seen through time reversal.  r28.1(a) found that the tachyon
measures ideally exactly when it couples through \(T\).  The two results are one statement: the coupling through \(T\) turns the
tachyon into the compensator partner.
\end{enumerate}

\paragraph{Proof sketch.}
(a) \((\tfrac12(1+J))^2=\tfrac14(1+2J+J^2)\).  (b) matrix product; \(z\mapsto i\bar z\) as in r02.  (c) combine (b) with r28.1(a).  \(\square\)

\subsection*{v4.9.0-r30 consequence for the reservoir}
\label{subsec:v490-r30-reservoir}

\paragraph{Theoremlet v4.9.0-r30.3.}
\begin{enumerate}[leftmargin=2em]
\item[(a)] If the exchange of r28 and r29 runs through the compensator, it is coherent by construction.  The Markov replacement of
r29.1(e) is the classical mixture that the compensator was introduced to avoid.
\item[(b)] With coherent exchange of strength \(\bsym^2\) the attractor is self-consistent exactly at the Hopf point (r28, r29).  The open
point of r29, why the exchange is coherent, is answered under the author's account of the compensator.
\item[(c)] At equal strength the coherent attractor is sharper.  At strength \(1/2\) after \(\pi/6\) rotations, coherent exchange gives
\(|b|=2/3\), port weights \(1/6\) and \(5/6\), entropy \(0.32\) bit.  Markov replacement gives \(|b|=0.577\), port weights \(1/4\) and \(3/4\),
entropy \(0.74\) bit.
\end{enumerate}

\paragraph{Proof sketch.}
(a) definition of the compensator, r30.1.  (b) r28.2, r29.1(e).  (c) iteration to the fixed point.  \(\square\)

\paragraph{On the wording of Theorem~\ref{thm:dark-energy}.}
That theorem calls \(\Omega_M=\bsym^2\) and \(\Omega_\Lambda=1-\bsym^2\) Markov transition probabilities.  Under the reading of this layer,
\(\bsym^2\) remains the weight of one step, \(\sin^2\theta\).  Successive steps are composed as amplitudes through the compensator, not as a
classical Markov chain.  The theorem keeps its text.  A status note under the revision rule of r25 would record this reading.  It is
not placed here, and the decision lies with the author.

\subsection*{v4.9.0-r30 status in the three frames}
\label{subsec:v490-r30-frames}

\begin{center}
\small
\begin{tabular}{@{}>{\raggedright\arraybackslash}p{0.30\textwidth}>{\raggedright\arraybackslash}p{0.62\textwidth}@{}}
\toprule
result & M / T / P\\
\midrule
r30.1 one half & the mixture of M with its T image gives \(1/2\); the coherent combination is definite at two phases; at P both give \(1/2\)\\
r30.2 partner through \(T\) & relates M and T; \(J_B\) keeps the side, \(z\mapsto iz\) changes it (r28)\\
r30.3 coherent reservoir & as r28.2(c) and r29; self-consistency at the Hopf point only\\
\bottomrule
\end{tabular}
\end{center}

\subsection*{v4.9.0-r30 register entries}
\label{subsec:v490-r30-register}

\begin{center}
\small
\begin{tabular}{@{}>{\raggedright\arraybackslash}p{0.34\textwidth}>{\raggedright\arraybackslash}p{0.58\textwidth}@{}}
\toprule
statement & status after v4.9.0-r30\\
\midrule
purpose of the compensator & to replace the classical one half by a definite port (author); confirmed by the main body's remark on light (r30.1)\\
compensator partner & the tachyon seen through time reversal, \(J_B=R_{\pi/2}\circ T\) (r30.2)\\
coherence of the exchange (r29) & answered: coherent by construction of the compensator (r30.3)\\
Markov wording of Thm.~\ref{thm:dark-energy} & read as one-step weight; status note possible, decision with the author\\
\bottomrule
\end{tabular}
\end{center}

\subsection*{v4.9.0-r30 ledger and guard update}
\label{subsec:v490-r30-ledger-guard}

\[
R_{\rm coh}^{v490\text{-}r30}=(0,0,0)
\]
for (one-half avoidance made precise, partner identified, coherence of the exchange answered).

The layer spends no guard:
\[
R_{\rm glob}^{(v490\text{-}r30,+)}=(0,1,1,1,1),
\qquad
\Gamma_{\rm guard}^{v490\text{-}r30}=(1,1,1,1,1),
\qquad
\beta_4=1/256\ \text{unspent}.
\]
It does not imply exact finite \(C5\) zero closure, exact all-readout terminality,
\(C6/T2\), a global photon-clock/RH theorem, or global all-readout closure.

